\chapter{Traditional User Pairing Methods}
\label{chap:traditional}

Having established the system model and channel generation framework in the previous chapter, we now turn our attention to traditional baseline user pairing strategies for NOMA systems. These methods represent the current state-of-the-art in practical NOMA deployments and optimization research, spanning from simple heuristic approaches requiring minimal computation to sophisticated graph-theoretic algorithms that guarantee globally optimal solutions. Understanding these traditional methods is essential for establishing performance benchmarks against which our proposed deep learning approach will be evaluated, and for appreciating the trade-offs between computational complexity and solution quality that motivate machine learning research in this domain.

We present four representative pairing strategies that collectively cover the design space from greedy heuristics through optimal combinatorial optimization. Static pairing employs the simplest possible strategy of matching strongest users with weakest users to maximize channel disparity. Balanced pairing moderates this approach by dividing users into upper and lower halves by channel quality and pairing within reasonable bounds. Blossom maximum-weight matching formulates the problem as general graph optimization and applies Edmonds' algorithm to find globally optimal solutions. Finally, bipartite graph matching with proportional fairness weighting provides a middle ground that exploits the natural strong-weak structure of NOMA while incorporating fairness objectives. For each method, we provide the algorithmic framework, computational complexity analysis, implementation details, and discuss advantages and limitations in the context of practical NOMA deployment.

\section{Static Pairing: Maximizing Channel Disparity}

Static pairing represents the most intuitive and computationally efficient approach to NOMA user clustering. The fundamental insight driving this method is that successive interference cancellation requires significant channel gain disparity between paired users, with typical implementations requiring at least 8 dB difference to ensure reliable SIC operation. By pairing users at opposite extremes of the channel quality distribution—matching the strongest user with the weakest, the second-strongest with the second-weakest, and so forth—static pairing guarantees maximum channel disparity and hence robust SIC feasibility.

The algorithm implementation is remarkably straightforward. First, all $N$ users are sorted in ascending order of channel power gain $|h_1|^2 \leq |h_2|^2 \leq \cdots \leq |h_N|^2$, which requires $O(N \log N)$ comparisons using efficient sorting algorithms such as quicksort or mergesort. Then, pairs are formed symmetrically from opposite ends of the sorted list: user 1 (weakest) is paired with user $N$ (strongest), user 2 with user $N-1$, and this process continues until all users are paired or one user remains unpaired if $N$ is odd. This pairing step requires only $O(N)$ operations, making the overall complexity $O(N \log N)$ which is dominated by the initial sorting phase.

The rationale for this approach extends beyond SIC feasibility considerations. By allocating the strongest possible partner to each weak user, static pairing ensures that cell-edge users with poor channels receive maximum assistance through high power allocation (since NOMA principles dictate that weaker users receive higher transmit power). This implicitly promotes fairness by preventing weak users from being paired with other weak users in which case neither could achieve acceptable rates. Additionally, the maximum channel disparity facilitates aggressive power allocation where a large fraction (often 70-80\%) of the PRB power can be allocated to the weak user while still allowing the strong user to achieve high rates after SIC due to its excellent channel quality.

From an implementation perspective, static pairing is extremely attractive for resource-constrained base stations or real-time applications with stringent latency requirements. The algorithm requires no complex optimization, no iterative search procedures, and minimal memory footprint—only storage for the sorted user list. It can be implemented in a few dozen lines of code and executes in microseconds even for thousands of users. These properties make static pairing a popular choice in early NOMA prototypes and testbeds where computational efficiency is paramount.

However, static pairing suffers from several fundamental limitations that constrain its performance in realistic deployment scenarios. Most significantly, it completely ignores spatial correlation and user geometry. Users at similar distances from the base station but with very different instantaneous channel realizations due to small-scale fading will be paired together, potentially violating the angular separation constraint we impose to ensure channel decorrelation. In scenarios with clustered user distributions, such as crowds in stadiums or public gatherings, static pairing may create geographically poor pairings that exacerbate interference rather than mitigating it.

Additionally, static pairing exhibits suboptimal performance for non-uniform user distributions. When channel gains are not uniformly spread across the dynamic range, pairing the extremes may sacrifice opportunities for better local pairings among users with moderate channel gains. For instance, if most users have similar weak channels with only a few strong users, static pairing will pair all weak users with the limited number of strong users, leaving many weak users unpaired rather than allowing weak-weak pairings that could still provide service albeit at reduced rates.

Finally, the fixed power allocation resulting from static pairing often wastes resources on strong users who could achieve near-capacity rates with much less power, while weak users might benefit from additional power that would increase their rates more substantially. The lack of joint optimization between pairing decisions and power allocation prevents static pairing from achieving global optimality in system sum-rate.

\section{Balanced Pairing: Moderating Disparity for Fairness}

Balanced pairing addresses some limitations of static pairing by moderating the channel gain disparity between paired users, thereby improving fairness while maintaining SIC feasibility. Rather than pairing extremes of the channel quality distribution, balanced pairing divides users into two groups—the weaker half and the stronger half—and pairs corresponding elements from each group. Specifically, after sorting users by channel gain as before, the list is split at the median: users $1$ through $N/2$ form the weak group, and users $(N/2 + 1)$ through $N$ form the strong group. User $i$ from the weak group is then paired with user $(N/2 + i)$ from the strong group.

This modification substantially changes the pairing structure. While static pairing creates a single pair with the maximum possible channel disparity (user 1 with user $N$), balanced pairing creates all pairs with approximately median disparity. The weakest user is now paired with the median user rather than the strongest, and the strongest user is paired with the user just below the median rather than the weakest. This reduces variance in per-pair channel disparity and creates more uniform rate distributions across users.

The computational complexity remains identical to static pairing at $O(N \log N)$ since the same sorting step is required, and the pairing phase still requires only $O(N)$ operations to form $N/2$ pairs. Thus, balanced pairing retains the attractive efficiency properties of static pairing while offering improved fairness characteristics.

The fairness improvement stems from more equitable power allocation patterns. In static pairing, the pair consisting of users 1 and $N$ receives vastly different power levels (perhaps 80\% to user 1 and 20\% to user $N$) due to extreme channel disparity, while in balanced pairing, all pairs receive more moderate power splits (perhaps 60\% to 40\%). This reduces the variance in individual user rates and improves Jain's fairness index, which penalizes high variance in rate distributions.

Additionally, balanced pairing produces shorter spatial connections on average. By pairing users with moderate channel differences rather than extreme differences, users that are geographically closer tend to be paired together, improving angular decorrelation and reducing the probability of violating spatial constraints. This property is particularly beneficial in clustered deployment scenarios where users naturally separate into weak and strong geographic regions.

However, balanced pairing inherits several limitations from static pairing while introducing new concerns. The fixed pairing rule still cannot adapt to channel statistics or exploit structure in specific scenarios—it applies the same median-split strategy regardless of whether the channel gain distribution is uniform, exponential, or has multiple modes. Furthermore, in some channel realizations, pairing users from the lower and upper halves may actually violate SIC constraints if the channel distribution is compressed such that the median-to-weakest disparity falls below the 8 dB threshold required for reliable SIC.

The method also fails to optimize for global metrics. While it achieves better fairness than static pairing, it does so by constraining the optimization space rather than through principled multi-objective optimization. There may exist alternative pairings that achieve even better fairness-throughput trade-offs, but balanced pairing cannot discover them due to its fixed rule. Despite these limitations, balanced pairing represents a pragmatic compromise between the simplicity of static pairing and the complexity of optimization-based approaches, making it a reasonable baseline for systems prioritizing fairness over absolute throughput maximization.

\section{Blossom Maximum-Weight Matching: Globally Optimal Solutions}

Moving beyond heuristic approaches, we now consider graph-theoretic optimization methods that formulate user pairing as a combinatorial matching problem on weighted graphs. The Blossom algorithm, developed by Jack Edmonds in 1965, represents one of the landmark achievements in combinatorial optimization, providing a polynomial-time algorithm for maximum-weight matching on general graphs—a problem that appears deceptively simple but requires sophisticated algorithmic techniques to solve efficiently.

To apply graph matching to NOMA user pairing, we construct an undirected graph $G = (V, E)$ where the vertex set $V$ corresponds to the $N$ users in the cell, and an edge $(i,j) \in E$ exists between users $i$ and $j$ if and only if pairing them satisfies all physical feasibility constraints. Specifically, edge $(i,j)$ is included in the graph if the channel gain disparity satisfies the SIC threshold—$10\log_{10}(|h_j|^2 / |h_i|^2) \geq 8$ dB assuming $|h_j|^2 > |h_i|^2$—and the angular separation constraint is satisfied—$|\theta_i - \theta_j| \geq 30°$ ensuring spatial decorrelation. This constraint enforcement during graph construction guarantees that any matching found by the algorithm will be physically feasible for NOMA operation.

Each edge $(i,j)$ is assigned a weight $w_{ij}$ representing the value or benefit of pairing users $i$ and $j$. For sum-rate maximization, we set $w_{ij} = R_{\text{sum}}(i,j,\alpha^*)$ where $\alpha^*$ is the optimal power allocation coefficient that maximizes the sum-rate for this specific pair. Computing $\alpha^*$ requires solving a one-dimensional optimization problem: $\alpha^* = \argmax_{\alpha \in (0,1)} [R_i(\alpha) + R_j(\alpha)]$ subject to SIC feasibility constraints. We employ golden section search or binary search to find $\alpha^*$ numerically, typically converging within 20-30 iterations to tolerance $10^{-4}$. This yields the maximum achievable sum-rate for pair $(i,j)$ which serves as the edge weight.

The matching problem is then to find a set of edges $\mathcal{M} \subseteq E$ such that each vertex appears in at most one edge (the matching constraint) and the total weight $\sum_{(i,j) \in \mathcal{M}} w_{ij}$ is maximized. This is precisely the maximum-weight matching problem on general graphs. Unlike bipartite matching which has more structure and simpler algorithms, general matching is considerably more complex due to the presence of odd cycles in the graph that can create ambiguities in augmenting path methods.

Edmonds' Blossom algorithm solves this problem through an elegant generalization of the Hungarian algorithm for bipartite graphs. The key innovation is the concept of "blossoms"—odd cycles in the graph that are contracted into supernodes during the algorithm execution, then expanded when an augmenting path passes through them. The algorithm maintains a dual solution providing lower bounds on the optimal matching value and iteratively improves both the primal matching and dual solution until optimality conditions are satisfied.

The worst-case computational complexity of the Blossom algorithm is $O(N^3)$ for $N$ vertices. This cubic scaling arises from the need to potentially explore $O(N)$ augmenting paths, each of which may require $O(N^2)$ operations to find and augment. In practice, efficient implementations using advanced data structures achieve better performance, and the algorithm often terminates in fewer than $N$ augmenting path iterations for sparse graphs. Nevertheless, the $O(N^3)$ complexity represents a significant computational burden compared to $O(N \log N)$ for heuristic methods, especially for large user populations.

We implement Blossom matching using the NetworkX library's \texttt{max\_weight\_matching} function, which provides a well-optimized implementation of Edmonds' algorithm. Our implementation first constructs the graph by enumerating all $\binom{N}{2}$ potential user pairs, checking feasibility constraints, computing optimal power allocation and sum-rates for feasible pairs, and adding edges to the graph with sum-rate as weight. This graph construction phase dominates the computation for small to moderate $N$, requiring $O(N^2)$ pairwise feasibility checks and power optimizations. The actual matching computation via Blossom then follows, returning a set of edges representing the optimal pairing.

The critical advantage of Blossom matching is global optimality. Unlike heuristics that make greedy or locally optimal decisions, Blossom is guaranteed to find the maximum-weight matching—no other pairing strategy can achieve higher total sum-rate given the same edge weights. This property makes Blossom matching an ideal benchmark for evaluating other algorithms. We use Blossom solutions as ground truth labels for training our Graph Neural Network, and as the performance upper bound against which all methods are compared.

However, Blossom matching faces limitations for practical deployment in large-scale networks. The $O(N^3)$ complexity becomes prohibitive as $N$ grows: for $N = 500$ users, our measurements show execution times approaching 45 seconds on commodity hardware, far exceeding the 1 millisecond TTI of 5G NR systems. This computational burden renders Blossom matching impractical for real-time resource allocation in dense deployments. Additionally, Blossom assumes that edge weights (sum-rates) are accurately computable a priori, requiring perfect CSI and the ability to solve the power allocation problem for each potential pair—assumptions that may not hold in dynamic channels with imperfect feedback.

\section{Bipartite Graph Matching with Proportional Fairness}

Bipartite graph matching represents a middle ground between heuristic methods and general graph optimization, exploiting the natural structure of NOMA user pairing to achieve near-optimal performance with improved computational efficiency compared to Blossom matching. The key observation is that NOMA inherently involves pairing weak users with strong users to enable SIC—weak-weak pairing provides no SIC opportunity, and strong-strong pairing wastes the potential of strong users who could successfully decode and cancel weak user interference. This suggests formulating the problem on a bipartite graph where one partition contains weak users and the other contains strong users.

We partition the user set by dividing the sorted channel gain list at the median: users $1$ through $N/2$ with channel gains below the median form the weak user set $V_1$, while users $(N/2 + 1)$ through $N$ with channel gains above the median form the strong user set $V_2$. The bipartite graph $G = (V_1 \cup V_2, E)$ contains edges only between $V_1$ and $V_2$—that is, edge $(i,j)$ exists only if $i \in V_1$, $j \in V_2$, and the pair satisfies SIC and angular separation constraints. This bipartite structure simplifies the matching problem considerably since odd cycles cannot exist in bipartite graphs, eliminating the need for blossom contraction.

An important enhancement we introduce is proportional fairness weighting rather than raw sum-rate edge weights. While maximizing total sum-rate is a natural objective, it can lead to highly unfair resource allocations where a few high-rate pairs monopolize resources at the expense of many low-rate users. Proportional fairness provides a principled framework for balancing efficiency and equity through logarithmic utility functions.

For each feasible edge $(i,j)$, we compute the power-optimized sum-rate $R_{\text{sum}}(i,j,\alpha^*)$ as before, but instead of using this directly as the edge weight, we apply a logarithmic transformation: $w_{ij}^{\text{PF}} = \log(R_i(\alpha^*) + \epsilon) + \log(R_j(\alpha^*) + \epsilon)$ where $\epsilon = 10^{-12}$ is a small constant preventing logarithm of zero. This transformation has several important properties. First, it is concave, meaning diminishing marginal utility—doubling a user's rate from 1 Mbps to 2 Mbps has greater impact than doubling from 100 Mbps to 200 Mbps. Second, the logarithm asymptotically approaches negative infinity as rate approaches zero, heavily penalizing zero-rate users and ensuring the objective prefers allocations where all users receive non-zero service. Third, maximizing sum of logarithms is equivalent to maximizing the geometric mean of rates, which is a classic fairness metric.

The bipartite maximum-weight matching problem can be solved using the Hungarian algorithm, also known as the Kuhn-Munkres algorithm, which was developed by Harold Kuhn in 1955 based on earlier work by König and Egerváry. The Hungarian algorithm solves the assignment problem—finding a perfect matching (all vertices matched) that maximizes total weight—in $O(N^3)$ time for an $N \times N$ bipartite graph. While this is the same asymptotic complexity as Blossom for general graphs, the Hungarian algorithm has better practical constants and is simpler to implement due to the bipartite structure.

Modern implementations further optimize performance through techniques such as sparse matrix representations when the bipartite graph has few edges, auction algorithms that achieve $O(NM \log N)$ complexity where $M$ is the number of edges, and parallel algorithms that distribute computation across multiple processors. We use SciPy's \texttt{linear\_sum\_assignment} function which implements an efficient variant of the Hungarian algorithm with numerical stability enhancements for floating-point weights.

Our implementation proceeds in several stages. First, we partition users into weak and strong sets based on the median channel gain. Second, we construct the bipartite graph by checking all $|V_1| \times |V_2|$ potential cross-partition pairs for feasibility constraints—SIC threshold and angular separation. Third, for each feasible pair, we optimize power allocation to maximize sum-rate and compute the proportional fairness weight using the logarithmic transformation. Fourth, we formulate the assignment problem as a cost matrix where entry $(i,j)$ contains the negated PF weight (negated because standard solvers minimize cost rather than maximize weight). Fifth, we invoke the Hungarian algorithm to find the optimal assignment. Finally, we extract the matched pairs and verify that all constraints remain satisfied.

Bipartite PF matching achieves several desirable properties. It naturally aligns with the weak-strong structure of NOMA, ensuring all pairs have sufficient channel disparity for SIC. The proportional fairness weighting balances throughput maximization with equity concerns, typically improving Jain's fairness index by 10-15\% compared to raw sum-rate optimization while sacrificing only 3-5\% of total throughput. The bipartite formulation reduces computational overhead compared to general graph matching, often achieving 20-30\% speedup compared to Blossom for the same problem size. The method also provides flexibility through the ability to adjust the weak-strong partition threshold: using the 40th percentile instead of median creates more weak users and fewer strong users, potentially improving coverage of cell-edge users at the expense of overall throughput.

Limitations include the somewhat arbitrary partition threshold which may not be optimal for all channel distributions, the restriction to perfect matchings which prevents leaving users unpaired even if no good partner exists, and the continued $O(N^3)$ scaling which remains problematic for very large user populations. Nevertheless, bipartite PF matching represents a strong baseline that achieves excellent performance-complexity trade-offs and serves as a natural comparison point for learning-based approaches.

\section{Computational Complexity Analysis and Comparison}

To provide concrete understanding of the computational burden imposed by each algorithm, we analyze both theoretical complexity and empirical runtime measurements. Table~\ref{tab:complexity_comparison} summarizes the key complexity characteristics.

\begin{table}[h]
\centering
\caption{Computational Complexity Comparison of Traditional Pairing Methods}
\label{tab:complexity_comparison}
\begin{tabular}{lccc}
\toprule
\textbf{Method} & \textbf{Time Complexity} & \textbf{Space Complexity} & \textbf{Optimality} \\
\midrule
Static Pairing & $O(N \log N)$ & $O(N)$ & Heuristic \\
Balanced Pairing & $O(N \log N)$ & $O(N)$ & Heuristic \\
Blossom Matching & $O(N^3)$ & $O(N^2)$ & Globally Optimal \\
Bipartite PF Matching & $O(N^3)$ & $O(N^2)$ & Optimal (bipartite) \\
\bottomrule
\end{tabular}
\end{table}

Static and balanced pairing achieve linearithmic $O(N \log N)$ complexity dominated by the sorting operation, making them suitable for real-time operation even with thousands of users. Their space complexity is linear $O(N)$, requiring storage only for user channel gains and pairing results.

Both graph-based methods exhibit cubic $O(N^3)$ time complexity, though this arises from different sources. For Blossom, the cubic term comes from augmenting path iterations in the matching algorithm itself. For bipartite matching, it comes from the Hungarian algorithm's dual update iterations. The space complexity is quadratic $O(N^2)$ due to explicit or implicit storage of the edge weight matrix.

The practical implications of these complexity differences are substantial. For $N = 12$ users, all methods execute in milliseconds, making complexity differences negligible. For $N = 100$ users, heuristics still execute in microseconds while graph methods require hundreds of milliseconds—noticeable but often acceptable. For $N = 500$ users representing dense deployments, heuristics execute in under 1 millisecond while Blossom requires 40-50 seconds, a difference of 5 orders of magnitude that clearly prohibits real-time Blossom operation.

This analysis reveals the fundamental motivation for machine learning approaches: if we can train a model offline using Blossom-generated optimal solutions as supervision, then deploy the trained model for real-time inference with complexity comparable to heuristics, we could achieve near-optimal performance with practical computational requirements. This is precisely the goal of our proposed Graph Neural Network framework presented in the following chapter.

\section{Summary and Insights}

This chapter has presented four traditional user pairing strategies representing the current state-of-the-art in NOMA research and practice. Static pairing offers maximum simplicity and efficiency but sacrifices optimization. Balanced pairing improves fairness while retaining efficiency. Blossom matching achieves global optimality at the cost of cubic complexity. Bipartite PF matching balances throughput, fairness, and computational cost through structured problem formulation.

Several key insights emerge from this analysis. First, a fundamental trade-off exists between solution quality and computational complexity—achieving theoretical optimality requires accepting impractical computation times for large-scale systems. Second, incorporating domain structure (the weak-strong nature of NOMA) into the problem formulation yields substantial benefits, as evidenced by bipartite matching's competitive performance. Third, fairness and efficiency are competing objectives that require explicit multi-objective optimization rather than hoping heuristics naturally balance them.

These traditional methods provide essential baselines for evaluating our proposed deep learning approach. Static and balanced pairing represent the lower complexity bound and heuristic performance levels we must exceed to justify machine learning's added sophistication. Blossom matching represents the upper performance bound we aspire to approach. Bipartite PF matching represents the current best practice balancing multiple objectives. In the next chapter, we present our Graph Neural Network framework designed to achieve Blossom-level performance with heuristic-level complexity, thereby bridging this critical gap for practical NOMA deployment.
